\documentclass[11pt,a4paper]{article}
\usepackage[utf8]{inputenc}
\usepackage[spanish]{babel}
\usepackage{geometry}
\geometry{top=2.5cm, bottom=2.5cm, left=2.5cm, right=2.5cm}
\usepackage{graphicx}
\usepackage{booktabs}
\usepackage{tabularx}
\usepackage{hyperref}
\usepackage{amsmath}
\usepackage{amssymb}
\usepackage{cite}

\hypersetup{
    colorlinks=true,
    linkcolor=blue,
    filecolor=magenta,      
    urlcolor=blue,
    citecolor=blue
}

\begin{document}

% --- CARÁTULA / ENCABEZADO TÉCNICO ---
\begin{center}
    {\Large \textbf{UNIVERSIDAD TÉCNICA ESTATAL DE QUEVEDO}} \\[0.2cm]
    {\large \textbf{FACULTAD DE CIENCIAS DE LA COMPUTACIÓN}} \\[0.2cm]
    {\large \textbf{CARRERA DE INGENIERÍA EN SOFTWARE}} \\[0.5cm]
    
    \vspace{0.3cm}
    {\Large \textbf{TA-IND-01: IEEE 29148 y Tipos de Requisitos de un Sistema Conocido}} \\[0.2cm]
    {\large Asignatura: Ingeniería de Requerimientos (ISR-401) · Paralelo A} \\[0.2cm]
    {\large Docente: Ing. Gleiston Cicerón Guerrero Ulloa, PhD} \\[0.4cm]
    
    \textbf{Estudiante:} Contreras Chávez Kevin Germán (1251167720) \\
    \textbf{Correo institucional:} \texttt{kcontrerasc3@uteq.edu.ec} \\
    \textbf{Repositorio:} \url{https://github.com/kcontrerasc3-cloud/ISR401-KevinContreras} \\
    \textbf{Fecha de entrega:} 7 de noviembre de 2026
\end{center}

\rule{\textwidth}{0.4pt}
\vspace{0.4cm}

% --- 1. DESCRIPCIÓN DEL SISTEMA ---
\section{Descripción del Sistema Seleccionado (Telegram)}
Para el desarrollo de esta actividad se ha seleccionado la plataforma de mensajería instantánea \textbf{Telegram}. Su propósito fundamental es proveer un servicio de comunicación multiplataforma basado en la nube, optimizado para ofrecer alta velocidad, sincronización multi-dispositivo en tiempo real y robustas capacidades para la gestión de comunidades masivas mediante canales, grupos y la integración de bots automatizados. Los usuarios principales van desde particulares que buscan mensajería rápida hasta administradores de comunidades y desarrolladores. El sistema opera en entornos conectados a internet mediante redes móviles (4G/5G) y Wi-Fi, ejecutándose de forma concurrente en dispositivos móviles (Android, iOS) y equipos de escritorio.

Las evidencias visuales de la plataforma se respaldan con tres capturas propias almacenadas en la ruta \texttt{TA-IND-01/figuras/}: el canal privado de pruebas configurado (\texttt{fig1\_canal.png}), la interfaz de ajustes de seguridad (\texttt{fig2\_seguridad.png}), y el panel de gestión de almacenamiento y red (\texttt{fig3\_almacenamiento.png}).

% --- 2. RESUMEN ANALÍTICO ---
\section{Resumen Analítico: Norma ISO/IEC/IEEE 29148:2018}
El proceso de ingeniería de requerimientos requiere la construcción de especificaciones rigurosas y libres de ambigüedad. De acuerdo con la norma \textbf{ISO/IEC/IEEE 29148:2018} (Cláusula 5.2.4), un \textbf{requisito bien formado} es aquel que cumple con una estructura sintáctica definida y propiedades lógicas que permiten su comprensión unívoca, su trazabilidad y su posterior verificación a lo largo del ciclo de vida del software. 

Asimismo, la cláusula 5.2.5 establece las características fundamentales que debe poseer un \textbf{requisito individual}, entre las cuales destacan: ser necesario (aporta valor directo al sistema), singular (contiene una sola idea o capacidad por sentencia), realizable (técnica y económicamente viable), verificable (existen medios empíricos o formales para comprobar su cumplimiento), correcto, inequívoco, modificable y trazable. Por otro lado, la cláusula 5.2.6 aborda las características de un \textbf{conjunto de requisitos}, exigiendo que este sea completo, consistente (sin contradicciones internas), no redundante y apto para ser validado por los interesados.

Finalmente, la cláusula 5.2.7 prohíbe explícitamente el uso de \textbf{términos ambiguos, vagos o subjetivos} (tales como ``user-friendly'', ``rápido'', ``aproximadamente'', ``robusto'', ``fácil'', ``flexible'' o conectores condicionales débiles como ``y/o''), ya que introducen múltiples interpretaciones e impiden establecer criterios objetivos de aceptación.

% --- 3. GLOSARIO ---
\section{Glosario (Cláusula 3)}
Con base en la cláusula 3 de la norma ISO/IEC/IEEE 29148:2018, se definen diez términos clave adaptados al contexto de Telegram:
\begin{enumerate}
    \item \textbf{Adquisición de requisitos (Elicitation):} Proceso de recolección y descubrimientos de necesidades con los usuarios. \textit{Ejemplo en Telegram:} Consulta de requerimientos para el envío de archivos multimedia pesados en chats privados.
    \item \textbf{Atributo de requisito:} Propiedad o metadato que caracteriza al requisito. \textit{Ejemplo en Telegram:} Asignar prioridad alta y trazabilidad al requisito de cifrado en chats secretos.
    \item \textbf{Línea base (Baseline):} Versión formalmente aprobada de un conjunto de requisitos que sirve como referencia para cambios futuros. \textit{Ejemplo en Telegram:} Versión 1.0 de las especificaciones de canales privados aprobada para desarrollo.
    \item \textbf{Parte interesada (Stakeholder):} Individuo u organización que afecta o es afectado por el sistema. \textit{Ejemplo en Telegram:} Administradores de canales y usuarios finales.
    \item \textbf{Requisito funcional:} Declaración de lo que el sistema debe hacer. \textit{Ejemplo en Telegram:} El sistema debe permitir la creación de canales privados con un enlace de invitación único.
    \item \textbf{Requisito no funcional (Calidad):} Criterio que delimita atributos del sistema. \textit{Ejemplo en Telegram:} Disponibilidad del servicio de almacenamiento en la nube.
    \item \textbf{Restricción:} Condición limitante impuesta al diseño o implementación. \textit{Ejemplo en Telegram:} Compatibilidad obligatoria con el protocolo MTProto para la transmisión de datos.
    \item \textbf{Trazabilidad (Traceability):} Habilidad de seguir la vida de un requisito en ambas direcciones. \textit{Ejemplo en Telegram:} Conectar el requisito de seguridad con su prueba de verificación en la interfaz de bloqueo.
    \item \textbf{Validación:} Confirmación mediante examen de que se cumplen los requisitos especificados para un uso previsto. \textit{Ejemplo en Telegram:} Pruebas de aceptación en la verificación en dos pasos.
    \item \textbf{Verificación:} Proceso de evaluación para determinar si los productos de una fase satisfacen las condiciones. \textit{Ejemplo en Telegram:} Comprobación técnica de que la latencia de entrega cumple el umbral estipulado.
\end{enumerate}

% --- 4. FORMULACIÓN DE REQUISITOS ---
\section{Formulación de 15 Requisitos de Telegram}
Se presentan quince requisitos del sistema distribuidos en cinco tipos de la norma, incorporando cinco requisitos de calidad evaluados bajo el modelo \textbf{ISO/IEC 25010:2023}.

\begin{itemize}
    \item \textbf{Tipo Funcional (RF):}
    \begin{enumerate}
        \item \textbf{REQ-FUN-01:} El sistema debe permitir la creación de un canal privado de difusión mediante la selección de opciones en el menú de configuración.
        \item \textbf{REQ-FUN-02:} El sistema debe permitir el envío de archivos multimedia de hasta 2 GB en chats individuales y canales.
        \item \textbf{REQ-FUN-03:} El sistema debe permitir al usuario activar la verificación en dos pasos ingresando una contraseña alfanumérica de respaldo.
        \item \textbf{REQ-FUN-04:} El sistema debe permitir la edición de mensajes enviados en un plazo máximo de 48 horas posteriores a su emisión.
        \item \textbf{REQ-FUN-05:} El sistema debe permitir la búsqueda de mensajes dentro de un chat mediante palabras clave introducidas en la barra de búsqueda superior.
    \end{enumerate}
    \item \textbf{Tipo de Calidad / Desempeño (RCA - ISO/IEC 25010: Eficiencia del desempeño):}
    \begin{enumerate}
        \item[6.] \textbf{REQ-CAL-01 (Desempeño):} El sistema debe entregar los mensajes de texto en un tiempo inferior a 300 ms bajo condiciones óptimas de red 4G. (\textit{Métrica:} Tiempo en milisegundos desde el clic de envío hasta la recepción; \textit{Umbral:} $\le 300$ ms).
        \item[7.] \textbf{REQ-CAL-02 (Fiabilidad - Madurez):} El servicio de infraestructura en la nube debe garantizar una disponibilidad operativa del 99.9\% anual. (\textit{Métrica:} Porcentaje de tiempo de actividad del servidor; \textit{Umbral:} $\ge 99.9\%$).
        \item[8.] \textbf{REQ-CAL-03 (Seguridad - Confidencialidad):} El sistema debe cifrar de extremo a extremo los mensajes en chats secretos utilizando el protocolo de cifrado MTProto 2.0. (\textit{Métrica:} Porcentaje de paquetes interceptados descifrables; \textit{Umbral:} $0\%$).
        \item[9.] \textbf{REQ-CAL-04 (Compatibilidad - Interoperabilidad):} El sistema debe sincronizar el historial de mensajes de forma concurrente en al menos 3 dispositivos distintos (móvil y escritorio) con una divergencia menor a 2 segundos. (\textit{Métrica:} Retraso de sincronización en segundos; \textit{Umbral:} $\le 2$ s).
        \item[10.] \textbf{REQ-CAL-05 (Usabilidad - Accesibilidad):} El sistema debe permitir la navegación completa de su interfaz mediante lectores de pantalla compatibles con sistemas operativos móviles. (\textit{Métrica:} Porcentaje de etiquetas de accesibilidad operativas; \textit{Umbral:} $100\%$).
    \end{enumerate}
    \item \textbf{Tipo de Interfaz Externa (RIE):}
    \begin{enumerate}
        \item[11.] \textbf{REQ-INT-01:} El sistema debe mostrar la interfaz gráfica de usuario adaptada al modo oscuro del sistema operativo del dispositivo.
        \item[12.] \textbf{REQ-INT-02:} El sistema debe conectarse a la API oficial de Telegram para la recepción y envío de datos mediante solicitudes HTTPS seguras.
    \end{enumerate}
    \item \textbf{Tipo de Restricción (RRES):}
    \begin{enumerate}
        \item[13.] \textbf{REQ-RES-01:} El código fuente del cliente oficial de escritorio debe compilarse utilizando bibliotecas de código abierto compatibles con arquitecturas de 64 bits.
    \end{enumerate}
    \item \textbf{Tipo de Proceso / Mantenibilidad (RPR - ISO/IEC 25010: Mantenibilidad):}
    \begin{enumerate}
        \item[14.] \textbf{REQ-MAN-01:} El sistema debe registrar un archivo de eventos de errores críticos (\textit{log}) en la memoria local del dispositivo ante fallos de sincronización. (\textit{Métrica:} Porcentaje de excepciones capturadas en logs; \textit{Umbral:} $100\%$).
        \item[15.] \textbf{REQ-MAN-02:} Las actualizaciones del paquete de instalación en dispositivos móviles no deben superar un tamaño de descarga comprimido de 150 MB. (\textit{Métrica:} Tamaño del archivo binario en megabytes; \textit{Umbral:} $\le 150$ MB).
    \end{enumerate}
\end{itemize}

% --- 5. TABLA DE CLASIFICACIÓN (USANDO TABULARX PARA CERO ADVERTENCIAS) ---
\section{Tabla de Clasificación y Criterios de Aceptación}
\begin{table}[h]
\centering
\scriptsize
\setlength{\tabcolsep}{3pt}
\begin{tabularx}{\textwidth}{l X l l X X}
\toprule
\textbf{ID} & \textbf{Enunciado Resumido} & \textbf{Tipo} & \textbf{ISO/IEC 25010} & \textbf{Métrica} & \textbf{Criterio de Aceptación} \\
\midrule
REQ-01 & Creación de canal privado & Funcional & N/A & Tasa de éxito & Canal creado y visible solo por invitación. \\
REQ-06 & Entrega rápida de mensajes & Calidad & Eficiencia (Desempeño) & Latencia (ms) & 100 de 100 mensajes entregados en $\le 300$ ms. \\
REQ-07 & Disponibilidad del servicio & Calidad & Fiabilidad (Madurez) & Porcentaje uptime & Disponibilidad anual $\ge 99.9\%$. \\
REQ-08 & Cifrado de extremo a extremo & Calidad & Seguridad (Confidencialidad) & Paquetes seguros & Cifrado MTProto activo sin fugas de texto plano. \\
REQ-09 & Sincronización multi-dispositivo & Calidad & Compatibilidad (Interoperabilidad) & Retraso (s) & Mensajes reflejados en $\le 2$ s en escritorio. \\
REQ-10 & Navegación accesible & Calidad & Usabilidad (Accesibilidad) & Cobertura UI & Lectores de pantalla leen el 100\% de menús. \\
\bottomrule
\end{tabularx}%
\caption{Matriz de clasificación y aceptación de requisitos seleccionados.}
\label{tab:clasificacion}
\end{table}

% --- 6. AUTOEVALUACIÓN Y CORRECCIÓN ---
\section{Autoevaluación de Requisitos}
\begin{table}[h]
\centering
\small
\begin{tabularx}{\textwidth}{l X X}
\toprule
\textbf{ID} & \textbf{Versión Anterior (Con fallos / términos vagos)} & \textbf{Versión Corregida (Norma ISO/IEC/IEEE 29148)} \\
\midrule
REQ-06 & El sistema debe enviar mensajes muy rápido a los chats. & El sistema debe entregar los mensajes de texto en un tiempo inferior a 300 ms bajo red 4G. \\
REQ-07 & El servidor de Telegram debe ser muy confiable siempre. & El servicio de infraestructura en la nube debe garantizar una disponibilidad anual del 99.9\%. \\
REQ-09 & La aplicación debe sincronizar bien y rápido en varias pantallas. & El sistema debe sincronizar el historial en al menos 3 dispositivos con divergencia $\le 2$ s. \\
\bottomrule
\end{tabularx}%
\caption{Tabla de autoevaluación y corrección de enunciados ambiguos.}
\label{tab:autoeval}
\end{table}

% --- 7. REFLEXIÓN PERSONAL ---
\section{Reflexión Personal}
Durante el desarrollo de esta actividad aplicada a la plataforma de Telegram, el aspecto que me representó mayor dificultad y requirió un análisis más profundo fue la formulación precisa de los requisitos de calidad cuantificables basados en el estándar ISO/IEC 25010:2023, específicamente al definir métricas y umbrales exactos para características no funcionales como el rendimiento y la fiabilidad. 

En la práctica habitual, uno suele describir las aplicaciones desde una perspectiva puramente funcional (lo que el sistema hace, como enviar mensajes, audios o archivos y crear canales), utilizando frases ambiguas o generales como «el sistema debe ser rápido» o «debe garantizar una buena seguridad». Sin embargo, al aplicar los criterios estipulados en la norma ISO/IEC/IEEE 29148:2018, comprendí que un requisito técnicamente válido exige eliminar toda ambigüedad. Esto implicó investigar y acotar parámetros numéricos reales, tales como establecer un umbral exacto de latencia en milisegundos para la entrega de mensajes o un porcentaje específico de disponibilidad del servicio en la nube. 

Este ejercicio me permitió valorar la importancia de la ingeniería de requisitos en el desarrollo de software, ya que pasar de enunciados vagos a especificaciones verificables es fundamental para evitar malentendidos entre desarrolladores y clientes, asegurando que el producto final sea medible y cumpla con altos estándares de calidad antes de su implementación.

% --- 8. REFERENCIAS BIBLIOGRÁFICAS ---
\section{Referencias Bibliográficas}
\begin{thebibliography}{4}
\bibitem{iso29148} ISO/IEC/IEEE, \emph{ISO/IEC/IEEE 29148:2018 Systems and software engineering -- Life cycle processes -- Requirements engineering}, International Standard, 2018.
\bibitem{iso25010} ISO/IEC, \emph{ISO/IEC/IEEE 25010:2023 Systems and software engineering -- Systems and software Quality Requirements and Evaluation (SQuaRE) -- Product quality model}, International Standard, 2023.
\bibitem{swebok} H. Washizaki, Ed., \emph{Guide to the Software Engineering Body of Knowledge (SWEBOK Guide)}, Version 4.0, IEEE Computer Society, 2024.
\end{thebibliography}

\end{document}